% !TEX root = ../Thesis.tex

\chapter{Implementierung und Ergebnisse}
\index{Implementierung}%
\index{Benchmark}%
\index{Evaluation}%
\label{ch:realisierung}

\section{Einleitung}

Ein voll besetzter Hörsaal. Hunderte Studierende greifen gleichzeitig zu ihren Smartphones und senden Antworten auf eine interaktive Frage über die Plattform \emph{frag.jetzt}. Innerhalb weniger Sekunden treffen zahlreiche Anfragen im Backend ein, während parallel Datenbankzugriffe erfolgen und externe Dienste angesprochen werden. Für die Nutzenden ist lediglich das Ergebnis sichtbar: Die Antwort erscheint auf der Leinwand – oder sie erscheint nicht.

Diese Situation verdeutlicht eine zentrale Herausforderung beim Betrieb moderner verteilter Webanwendungen. Wenn eine Anfrage ungewöhnlich lange benötigt oder vollständig fehlschlägt, stellt sich unmittelbar die Frage nach der Ursache. Liegt das Problem im Frontend, im Backend, in der Datenbank oder in einem externen Dienst? In verteilten Architekturen ist die Diagnose solcher Probleme ohne geeignete Betriebsdaten erheblich erschwert, da einzelne Systemkomponenten unabhängig voneinander arbeiten und Fehlerzustände nicht geradewegs sichtbar sind (Usman et al., 2022).

\subsection{Ausgangssituation}

Die Plattform \emph{frag.jetzt} ist eine Open-Source-Frage- und Antwortplattform, die insbesondere in Lehrveranstaltungen eingesetzt wird, um interaktive Rückmeldungen von Studierenden zu sammeln (TH Mittelhessen \& arsnova, 2025). Technisch basiert die Anwendung auf einer verteilten Architektur mit einem Progressive-Web-App-Frontend, einem Spring-Boot-Backend, einer PostgreSQL-Datenbank sowie externen Diensten, beispielsweise für KI-gestützte Funktionen.

Solche Architekturen bieten eine hohe Flexibilität sowie eine gute Skalierbarkeit, erhöhen jedoch gleichzeitig die Komplexität des Betriebs. Während moderne Entwicklungspraktiken wie Continuous Integration und Continuous Delivery schnelle und automatisierte Softwarebereitstellungen ermöglichen, bleibt häufig unklar, wie sich das Verhalten der Anwendung im realen Betrieb nachvollziehen lässt.

Traditionelle Monitoring-Systeme erfassen dazu typischerweise quantitative Betriebsdaten wie Anfragevolumen, Fehlerraten oder Antwortzeiten. Monitoring bezeichnet dabei die kontinuierliche Sammlung, Aggregation und Darstellung von Systemmetriken zur Bewertung des Systemzustands (Ewaschuk, 2017). In stark verteilten Systemen reichen solche Ansätze jedoch häufig nicht aus, um komplexe Fehlerzustände oder Leistungsprobleme eindeutig zu analysieren.

Observability erweitert diesen Ansatz, indem sie es ermöglicht, aus verschiedenen Telemetriedaten Rückschlüsse auf den internen Zustand eines Systems zu ziehen. Typische Datenquellen sind \emph{Metriken}, \emph{Logs} sowie \emph{Distributed Traces} (Faseeha et al., 2025). Erst durch die Kombination dieser Daten entsteht eine ganzheitliche Sicht auf das Verhalten eines Systems und seine internen Abhängigkeiten.

\subsection{Problemstellung}

Trotz der technischen Reife moderner Entwicklungsprozesse fehlt in vielen Anwendungen eine konsistente Strategie zur Beobachtung des Systembetriebs. Dies gilt auch für Plattformen wie \emph{frag.jetzt}, deren Architektur aus mehreren voneinander abhängigen Komponenten besteht. Ohne klar definierte Metriken sowie ohne korrelierte Telemetriedaten bleibt der interne Zustand der Anwendung weitgehend intransparent.

Ein etabliertes Konzept zur Bewertung der Betriebsqualität verteilter Systeme sind die sogenannten \emph{Four Golden Signals}. Diese umfassen die vier Messgrößen \emph{Latency}, \emph{Traffic}, \emph{Errors} und \emph{Saturation} und ermöglichen eine kompakte Bewertung des Systemzustands anhand zentraler Leistungsindikatoren (Ewaschuk, 2017).

Ein weiteres Problem besteht in der fehlenden Korrelation unterschiedlicher Telemetriedaten. In Cloud-Native-Anwendungen entstehen Logeinträge, Metriken und Traces an verschiedenen Stellen der Systemarchitektur. Ohne geeignete Mechanismen zur Zusammenführung dieser Daten wird es schwierig, einzelne Anfragen über mehrere Komponenten hinweg nachzuverfolgen und mögliche Ursachen von Fehlern zu identifizieren (Albuquerque \& Correia, 2025).

Darüber hinaus kann eine undurchdachte Alarmierungsstrategie zu sogenannter \emph{Alert Fatigue} führen. Dabei werden so viele Warnmeldungen erzeugt, dass wichtige Hinweise im täglichen Betrieb übersehen oder ignoriert werden. Studien zeigen, dass eine gezielte Priorisierung von Alarmmeldungen entscheidend ist, um Betriebsteams bei der effektiven Reaktion auf Störungen zu unterstützen (Jalalvand et al., 2024).

Diese Herausforderungen verdeutlichen, dass nicht nur die Erhebung von Betriebsdaten erforderlich ist, sondern auch eine strukturierte Strategie zu deren Auswahl, Aggregation und Interpretation.

\subsection{Zielsetzung der Arbeit}

Absicht dieser Arbeit ist die Entwicklung einer umfassenden Observability-Strategie für die Plattform \emph{frag.jetzt}. Dabei wird ein Ansatz entwickelt, der relevante Telemetriedaten definiert, geeignete Observability-Werkzeuge identifiziert und ein effektives Alarmierungskonzept entwirft.

Die Strategie umfasst mehrere Ebenen der Systembeobachtung. Dazu gehören insbesondere Metriken zur Leistungsanalyse, Protokollierung, verteiltes Tracing sowie Konzepte für Dashboards und Alarmmeldungen. Durch diese Kombination soll es möglich werden, den Zustand der Anwendung kontinuierlich zu analysieren und Ursachen für Fehler oder Leistungsprobleme schneller zu identifizieren.

Neben der konzeptionellen Entwicklung steht auch die praktische Umsetzbarkeit im Fokus. Ziel ist daher eine Observability-Strategie, die sich in reale Entwicklungs- und Betriebsumgebungen integrieren lässt und konkrete Handlungsempfehlungen (Runbooks) für den Betrieb der Plattform liefert.

\subsection{Forschungsfragen}

Aus der beschriebenen Problemstellung ergeben sich die folgenden Forschungsfragen:

\begin{enumerate}
\item Wie können die Four Golden Signals für die spezifischen Services von frag.jetzt (PWA-Frontend, Spring Boot Backend, PostgreSQL, KI-Services) konkret definiert und instrumentiert werden?

\item Welche Observability-Stack-Kombination ist am besten geeignet für die Monitoring-Anforderungen von frag.jetzt?

\item Wie kann ein intelligentes Alert-System entwickelt werden, das False Positives minimiert und actionable Insights für verschiedene Stakeholder-Rollen (DevOps, Entwickler, Product Owner) bereitstellt?
\end{enumerate}

\subsection{Aufbau der Arbeit}

Die Arbeit ist in mehrere aufeinander aufbauende Kapitel gegliedert. Kapitel 2 erläutert zunächst die theoretischen Grundlagen von Monitoring und Observability und stellt zentrale Konzepte verteilter Systeme dar. Kapitel 3 analysiert anschließend die Architektur der Plattform \emph{frag.jetzt} und identifiziert observability-relevante Komponenten.

Kapitel 4 beschreibt das methodische Vorgehen zur Entwicklung der Observability-Strategie. Darauf aufbauend definiert Kapitel 5 einen Metrikkatalog auf Basis der Four Golden Signals für die einzelnen Services der Plattform.

Kapitel 6 untersucht verschiedene Observability-Stacks und bewertet deren Eignung für den Einsatz in der betrachteten Architektur. Kapitel 7 beschreibt die Implementierung der Telemetrie-Pipeline sowie die Instrumentierung der Anwendung.

Kapitel 8 widmet sich der Gestaltung rollenbasierter Dashboards. Kapitel 9 entwickelt ein Konzept für intelligentes Alerting sowie exemplarische Runbooks für typische Störungsszenarien.

Abschließend diskutiert Kapitel 10 die Ergebnisse und mögliche Limitationen der Arbeit, bevor Kapitel 11 die wichtigsten Erkenntnisse zusammenfasst und einen Ausblick auf zukünftige Entwicklungen gibt.



\section{Grundlagen und theoretische Fundierung}

Die Sicherstellung eines zuverlässigen Betriebs von Softwaresystemen erfordert eine methodische Herangehensweise. Insbesondere in stark verteilten Umgebungen reicht die alleinige Betrachtung isolierter Hardware-Ressourcen nicht mehr aus. Dieses Kapitel legt das theoretische Fundament für die spätere Systemanalyse. Es grenzt grundlegende Begrifflichkeiten voneinander ab. Anschließend werden die maßgeblichen Kategorien von Telemetriedaten beleuchtet, bevor spezifische Herausforderungen von Microservices sowie bewährte Metriken und Alarmierungsstrategien erörtert werden.

\subsection{Monitoring und Observability}

Die Begriffe Monitoring und Observability werden in der Literatur häufig synonym verwendet, obwohl sie unterschiedliche Konzepte beschreiben. Monitoring umfasst das fortlaufende Sammeln, Verarbeiten und Anzeigen quantitativer Daten eines Systems in Echtzeit (Ewaschuk, 2017, S. 2 [SEITENZAHL PRÜFEN]). Es dient primär dazu, vordefinierte Fehlzustände zu erkennen. Ein Monitoring-System beantwortet die Frage, ob ein System ordnungsgemäß funktioniert. Dies ist essenziell für den Basisbetrieb.

Observability hingegen entstammt ursprünglich der klassischen Kontrolltheorie. Sie beschreibt die Fähigkeit, den internen Zustand eines Systems allein anhand seiner externen Ausgaben nachvollziehen zu können (Abieba et al., 2025, S. 271 [SEITENZAHL PRÜFEN]). Während traditionelles Monitoring auf statische Dashboards und bekannte Fehler ("known unknowns") ausgelegt ist, zielt Observability auf das Verstehen unvorhersehbarer Verhaltensweisen ("unknown unknowns") ab (Sakinala, 2025, S. 102 [SEITENZAHL PRÜFEN]). Observability beantwortet somit die Frage, warum ein System nicht funktioniert. Für die Erreichung dieses Zustands ist eine tiefergehende Instrumentierung des Quellcodes unabdingbar.

[HIER MÖGLICHES BILD EINFÜGEN: Gegenüberstellung von Monitoring (Was ist kaputt?) und Observability (Warum ist es kaputt?) anhand eines Schichtenmodells]

\subsection{Telemetriedaten in Observability-Systemen}

Um einen hohen Grad an Observability zu erreichen, müssen Systeme detaillierte Telemetriedaten emittieren. Diese Daten werden in der Fachliteratur oft als die drei Säulen der Observability bezeichnet: Logs, Metrics und Distributed Tracing (Albuquerque \& Correia, 2025, S. 2 [SEITENZAHL PRÜFEN]).

Logs sind unveränderliche, zeitgestempelte Aufzeichnungen diskreter Ereignisse, die innerhalb einer Applikation auftreten. Sie enthalten detaillierte Kontextinformationen zu spezifischen Transaktionen oder Fehlermeldungen (Banerjee \& Singh, 2024, S. 2 [SEITENZAHL PRÜFEN]). Da unstrukturierte Text-Logs maschinell schwer durchsuchbar sind, wird in der Praxis zunehmend auf strukturierte Formate wie JSON gesetzt. Dies erleichtert die automatisierte Aggregation.

Metrics stellen numerische Messwerte dar. Diese werden über festgelegte Zeitintervalle hinweg aggregiert (Sakinala, 2025, S. 103 [SEITENZAHL PRÜFEN]). Typische Beispiele sind die CPU-Auslastung, der belegte Arbeitsspeicher oder die Anzahl fehlerhafter HTTP-Anfragen. Metriken sind äußerst speicherplatzeffizient. Sie eignen sich hervorragend für die Erstellung von Dashboards und die Auslösung von Alarmen.

Distributed Tracing ermöglicht die lückenlose Verfolgung einer Anfrage über verschiedene Service-Grenzen hinweg. Ein Trace repräsentiert den gesamten Ausführungspfad einer Nutzerinteraktion. Er besteht aus mehreren "Spans", wobei jeder Span eine einzelne Arbeitseinheit innerhalb eines bestimmten Services darstellt (Albuquerque \& Correia, 2025, S. 3 [SEITENZAHL PRÜFEN]). Durch die Injektion einer eindeutigen Trace-ID in die Header der Netzwerkkommunikation lassen sich Latenzengpässe in komplexen Architekturen präzise lokalisieren.

\subsection{Observability in Microservice-Architekturen}

Die Verschiebung von monolithischen Architekturen hin zu Microservices bietet enorme Vorteile hinsichtlich Skalierbarkeit und unabhängiger Deployments. Gleichzeitig steigt die operative Komplexität drastisch an (Faseeha et al., 2025, S. 72012 [SEITENZAHL PRÜFEN]). Ein einzelner Aufruf durch einen Endnutzer kann Dutzende nachgelagerte API-Kommunikationen und Datenbankzugriffe auslösen.

Eine der typischen Herausforderungen besteht in der Bewältigung verteilter Fehlerzustände. Wenn ein externer Service langsam antwortet, kann dies zu einem Rückstau von Anfragen führen. Solche kaskadierenden Fehler bringen im schlimmsten Fall das gesamte System zum Erliegen (Usman et al., 2022, S. 86906 [SEITENZAHL PRÜFEN]). Herkömmliche Überwachungswerkzeuge stoßen hier an ihre Grenzen, da sie oft nur die Metriken isolierter Container betrachten.

Es bedarf einer End-to-End Visibility. Die gesamte Kette vom initialen Klick im Frontend bis hin zur Datenbankabfrage im Backend muss transparent gemacht werden (Fischer et al., 2024, S. 205 [SEITENZAHL PRÜFEN]). Nur so können Entwickler und operative Teams im Fehlerfall zweifelsfrei feststellen, an welchem Knotenpunkt innerhalb des Netzwerks die Störung ihren Ursprung hat.

\subsection{Die Four Golden Signals}

Um die Informationsflut an Metriken handhabbar zu machen, haben sich in der Praxis standardisierte Indikatoren etabliert. Die von Google formulierten "Four Golden Signals" bilden ein bewährtes Rahmenwerk zur Überwachung benutzerorientierter Dienste (Ewaschuk, 2017, S. 3 [SEITENZAHL PRÜFEN]).

Erstens misst die Latency (Latenz) die Zeitdauer, die ein System für die Verarbeitung einer Anfrage benötigt. Hierbei ist es entscheidend, die Latenz von erfolgreichen und fehlerhaften Anfragen getrennt zu betrachten (Ewaschuk, 2017, S. 4 [SEITENZAHL PRÜFEN]). Zweitens beschreibt der Traffic (Verkehrsaufkommen) die Belastung des Systems. Je nach Natur des Dienstes wird dies in HTTP-Anfragen pro Sekunde oder der Anzahl aktiver Nutzersessions gemessen.

Drittens erfassen die Errors (Fehlerraten) den Anteil der Anfragen, die fehlschlagen. Dies schließt explizite HTTP-500-Fehler ebenso ein wie logische Fehler, bei denen zwar der korrekte Statuscode gesendet wird, die übertragenen Daten jedoch unvollständig sind (Sakinala, 2025, S. 105 [SEITENZAHL PRÜFEN]). Viertens gibt die Saturation (Sättigung) an, wie stark die Ressourcenkapazität des Systems ausgelastet ist. Da viele Systeme bereits vor Erreichen einer 100-prozentigen Auslastung Performance-Einbußen zeigen, dient dieser Indikator als Frühwarnsystem (Ewaschuk, 2017, S. 5 [SEITENZAHL PRÜFEN]).

In modernen, Cloud-nativen Architekturen werden diese Signale zunehmend adaptiert. Während klassische Ansätze physische Ressourcen im Blick hatten, fokussieren sich heutige Praktiken stark auf Service Level Objectives (SLOs) und die tatsächliche Nutzererfahrung an den Schnittstellen der Container (Faseeha et al., 2025, S. 72015 [SEITENZAHL PRÜFEN]).

[HIER MÖGLICHES BILD EINFÜGEN: Tabellarische Übersicht der Four Golden Signals mit Beispielen für Backend- und Datenbank-Metriken]

\subsection{Alerting und Incident Management}

Das Sammeln von Metriken stiftet nur dann Wert, wenn daraus rechtzeitig operative Handlungen abgeleitet werden. Fehlkonfigurierte Alarmierungssysteme führen jedoch häufig zum Phänomen der Alert Fatigue (Alarmmüdigkeit). Wenn Entwickler permanent mit unkritischen Benachrichtigungen überflutet werden, sinkt die Aufmerksamkeit. Echte Notfälle bleiben infolgedessen oft unbemerkt [NEUE QUELLE].

Um dieser Ermüdung entgegenzuwirken, ist eine strikte Priorisierung unerlässlich. Ein robuster Ansatz ist das SLO-basierte Alerting. Anstatt einen Alarm auszulösen, sobald die CPU-Auslastung eines Servers 80 Prozent überschreitet, wird das System so konfiguriert, dass es nur dann alarmiert, wenn die Fehlerrate für den Endnutzer spürbar steigt oder das zugewiesene Error Budget rapide abnimmt (Jalalvand et al., 2024, S. 5 [SEITENZAHL PRÜFEN]). Symptome, die den Nutzer direkt betreffen, erhalten höchste Priorität.

Ursachenorientierte Metriken dienen hingegen primär der Fehlersuche im Nachgang und sollten lediglich auf Dashboards sichtbar sein, anstatt aktive Alarme auf Mobiltelefonen auszulösen (Ewaschuk, 2017, S. 15 [SEITENZAHL PRÜFEN]). Diese Strategie schont die mentalen Ressourcen der Teams und fördert eine nachhaltige Incident-Management-Kultur.


\section{Grundlagen und theoretische Fundierung}

Die Sicherstellung eines zuverlässigen Betriebs von Softwaresystemen erfordert eine methodische Herangehensweise. Insbesondere in stark verteilten Umgebungen reicht die alleinige Betrachtung isolierter Hardware-Ressourcen nicht mehr aus, um den Gesundheitszustand der Anwendung verlässlich zu bewerten. Dieses Kapitel legt das theoretische Fundament für die anschließende Systemanalyse von frag.jetzt. Zu Beginn werden grundlegende Begrifflichkeiten voneinander abgegrenzt. Darauf aufbauend werden die maßgeblichen Kategorien von Telemetriedaten beleuchtet, bevor spezifische Herausforderungen verteilter Dienste sowie bewährte Indikatoren und Alarmierungsstrategien detailliert erörtert werden.

\subsection{Monitoring und Observability}

Obwohl die Begriffe Monitoring und Observability in der Fachliteratur oft synonym auftauchen, beschreiben sie grundverschiedene Konzepte. Monitoring umfasst das fortlaufende Sammeln, Verarbeiten, Aggregieren und Anzeigen quantitativer Daten eines Systems in Echtzeit (Ewaschuk, 2017, S. 1). Es dient primär der Erkennung vordefinierter Fehlzustände. Ein solches System beantwortet die Frage, ob eine Anwendung ordnungsgemäß funktioniert, was für den Basisbetrieb unerlässlich ist.

Observability hingegen beschreibt die Fähigkeit, den internen Zustand eines Systems allein anhand seiner externen Ausgaben nachvollziehen zu können (Abieba et al., 2025, S. 270). Während herkömmliche Überwachungsansätze auf vorhersehbare Fehlermuster ("known unknowns") ausgerichtet sind, zielt Observability auf das Verstehen unvorhersehbarer und neuartiger Verhaltensweisen ("unknown unknowns") ab (Banerjee, 2025, S. 916). Folglich beantwortet Observability die tiefergehende Frage, warum ein System nicht funktioniert. Für die Erreichung dieses Zustands ist eine umfassende Instrumentierung des Quellcodes unabdingbar.

\subsection{Telemetriedaten in Observability-Systemen}

Um ein hohes Maß an Observability zu erzielen, müssen Anwendungen detaillierte Telemetriedaten emittieren. Diese Daten werden in der Fachliteratur häufig als die drei Säulen der Observability bezeichnet: Logs, Metrics und Traces (Albuquerque \& Correia, 2025, S. 2).

Logs stellen unveränderliche, zeitgestempelte Aufzeichnungen diskreter Ereignisse dar. Sie liefern tiefe Einblicke und essenzielle Kontextinformationen zu spezifischen Transaktionen oder Fehlermeldungen innerhalb eines Systems (Banerjee, 2025, S. 917). Da reine Textprotokolle maschinell schwer durchsuchbar sind, wird in der Praxis zunehmend auf definierte Formate wie JSON gesetzt, um die automatisierte Aggregation auszuwerten.

Metrics sind numerische Messwerte, welche über festgelegte Zeitintervalle hinweg aggregiert werden (Sakinala, 2025, S. 102). Typische Beispiele umfassen die CPU-Auslastung, den belegten Arbeitsspeicher oder die Anzahl fehlerhafter HTTP-Anfragen. Metriken zeichnen sich durch eine hohe Speicherplatzeffizienz aus. Sie eignen sich hervorragend für die Erstellung von Dashboards und die Auslösung von Alarmen.

Distributed Tracing ermöglicht die lückenlose Verfolgung einer Anfrage über verschiedene Service-Grenzen hinweg. Ein Trace repräsentiert den gesamten Ausführungspfad einer Nutzerinteraktion. Er besteht aus mehreren Spans, wobei jeder Span eine einzelne Arbeitseinheit innerhalb eines bestimmten Services abbildet (Albuquerque \& Correia, 2025, S. 6). Durch die Injektion einer eindeutigen Trace-ID in die Header der Netzwerkkommunikation lassen sich Latenzengpässe in feingranularen Architekturen präzise lokalisieren.

\subsection{Observability in Microservice-Architekturen}

Die Abkehr von monolithischen Architekturen hin zu verteilten Microservices bietet enorme Vorteile hinsichtlich Skalierbarkeit und unabhängiger Deployments. Gleichzeitig steigt die operative Komplexität drastisch an (Faseeha et al., 2025, S. 72012). Ein einzelner Aufruf durch einen Endnutzer kann Dutzende nachgelagerte API-Kommunikationen und Datenbankzugriffe auslösen.

Eine der größten Herausforderungen besteht in der Bewältigung verteilter Fehlerzustände. Wenn ein externer Service langsam antwortet, kann dies zu einem Rückstau von Anfragen führen. Solche Ausbreitungsmuster verursachen kaskadierende Fehler, die sich innerhalb von Millisekunden über mehrere Service-Grenzen hinweg fortpflanzen können (Sakinala, 2025, S. 101). Herkömmliche Überwachungswerkzeuge stoßen in diesem Kontext an ihre Grenzen, da sie oft nur die Metriken isolierter Container betrachten.

Es bedarf einer lückenlosen Sichtbarkeit. Die gesamte Kette vom initialen Klick im Frontend bis hin zur Datenbankabfrage im Backend muss transparent gemacht werden (Sakinala, 2025, S. 103). Nur durch diese vollständige Nachverfolgbarkeit können Entwickler und operative Teams im Fehlerfall zweifelsfrei feststellen, an welchem Knotenpunkt innerhalb des Netzwerks die Störung ihren Ursprung hat.

\subsection{Die Four Golden Signals}

Um die Informationsflut an Metriken handhabbar zu gestalten, haben sich in der Praxis standardisierte Indikatoren etabliert. Die von Google formulierten "Four Golden Signals" bilden ein bewährtes Rahmenwerk zur Beurteilung benutzerorientierter Dienste (Ewaschuk, 2017, S. 7).

Erstens misst die Latency (Latenz) die Zeitdauer, die ein System für die Verarbeitung einer Anfrage benötigt. Hierbei ist es entscheidend, die Latenz von erfolgreichen und fehlerhaften Anfragen strikt getrennt zu betrachten. Zweitens beschreibt der Traffic (Verkehrsaufkommen) die Belastung der Anwendung. Je nach Natur des Dienstes wird dieser Wert in HTTP-Anfragen pro Sekunde oder der Anzahl aktiver Nutzersessions gemessen.

Drittens erfassen die Errors (Fehlerraten) den Anteil der Anfragen, die fehlschlagen. Dies schließt explizite HTTP-500-Fehler ebenso ein wie logische Fehler, bei denen zwar der korrekte Statuscode gesendet wird, die gelieferten Inhalte jedoch fehlerhaft sind (Ewaschuk, 2017, S. 7). Viertens gibt die Saturation (Sättigung) an, wie stark die Ressourcenkapazität des Systems ausgelastet ist. Da viele Dienste bereits vor Erreichen einer vollständigen Auslastung Performance-Einbußen zeigen, dient dieser Indikator als wichtiges Frühwarnsystem.

In modernen, Cloud-nativen Architekturen werden diese Signale zunehmend adaptiert. Während klassische Ansätze physische Hardwareressourcen im Blick hatten, fokussieren sich heutige Praktiken stark auf Service Level Objectives (SLOs) und die tatsächliche Nutzererfahrung an den Schnittstellen der Dienste (Faseeha et al., 2025, S. 72012).

\subsection{Alerting und Incident Management}

Das Sammeln von Metriken stiftet nur dann Mehrwert, wenn daraus rechtzeitig operative Handlungen abgeleitet werden. Fehlkonfigurierte Alarmierungssysteme führen jedoch häufig zur sogenannten Alert Fatigue (Alarmmüdigkeit). Wenn das Personal permanent mit unkritischen Benachrichtigungen überflutet wird, sinkt die Aufmerksamkeit, was unweigerlich dazu führt, dass echte Notfälle verzögert bearbeitet oder gänzlich übersehen werden (Jalalvand et al., 2024, S. 2).

Um dieser Ermüdung entgegenzuwirken, ist eine strikte Priorisierung unerlässlich. Ein robuster Ansatz ist die Ausrichtung an Service Level Objectives. Anstatt einen Alarm auszulösen, sobald die CPU-Auslastung eines Servers einen fixen Schwellenwert überschreitet, wird das System so konfiguriert, dass es nur dann alarmiert, wenn die Verfügbarkeit oder Latenz für den Endnutzer spürbar leidet (Sakinala, 2025, S. 103). Symptome, die den Nutzer direkt betreffen, erhalten bei der Priorisierung stets höchste Dringlichkeit (Ewaschuk, 2017, S. 7).

Ursachenorientierte Metriken dienen hingegen primär der tiefergehenden Fehlersuche im Nachgang. Diese Detailinformationen sollten übersichtlich auf Dashboards visualisiert werden, anstatt aktive Alarme auf den Mobiltelefonen der Entwickler auszulösen (Ewaschuk, 2017, S. 5). Eine derartige Strategie schont die mentalen Kapazitäten der Teams und fördert eine nachhaltige Incident-Management-Kultur.


\section{Implementierungsdetails}

Das Kapitel dokumentiert die praktischen Aspekte der Umsetzung\index{Implementierungspraxis}.\footnote{Implementierungsdetails orientieren sich an Best Practices aus~\cite{martin2008clean}.}
Die Implementierung erfolgte iterativ\index{Iterative Entwicklung} über einen Zeitraum von 4 Monaten mit kontinuierlichem Testen\index{Continuous Testing} und Verfeinerung.

\subsection{Technologiestack}

Folgende Technologien wurden für die Implementierung eingesetzt:

\begin{sloppypar}
\begin{description}
  \item[Programmiersprache:] Python 3.11+\index{Python} (Typ-Hinweise\index{Type Hints}, Async/Await-Un\-ter\-stüt\-zung\index{Async/Await})
  \item[LLM-Integration:] OpenAI API (GPT-4)\index{GPT-4}, Anthropic API (Claude 3.5 Sonnet)\index{Claude}, mit Fallback-Mechanismus\index{Fallback-Mechanismus}
  \item[Orchestrierung:] LangChain 0.1.x für Basis-Komponenten\index{LangChain}, benutzerdefinierte Erweiterungen für SE-spezifische Tools; strukturierte Tool-Calls über \textbf{JSON}-Schemas und Konfigurationen per \textbf{YAML}.
  \item[Gedächtnis:] Chroma\index{Chroma Vector-DB} Vektordatenbank für semantisches Retrieval, SQLite für episodische Logs\index{Event-Logging}
  \item[Werkzeuglaufzeitumgebung:] Docker 24.0+\index{Docker!Sandboxing} für Sandboxing, pytest für Tests, ruff/black für Linting\index{Code Linting}
  \item[Monitoring:] Prometheus für Metriken\index{Prometheus}, strukturierte Logs (JSON) mit Trace-IDs\index{Distributed Tracing}
\end{description}
\end{sloppypar}

\subsection{Projektstruktur}

Das Projekt folgt einer modularen Architektur:

\begin{lstlisting}[breaklines=true, basicstyle=\ttfamily\small]
agent_se/
+-- core/
|   +-- agent.py           # Agent Controller (ReAct-Loop)
|   +-- policy.py          # Planning & Reflection Logic
|   +-- state.py           # State Management
+-- tools/
|   +-- base.py            # Tool Interface & Registry
|   +-- code_tools.py      # Linter, Formatter, AST-Parser
|   +-- test_tools.py      # pytest, coverage, Test-Runner
|   +-- vcs_tools.py       # git operations
+-- memory/
|   +-- episodic.py        # Event Logging & Retrieval
|   +-- semantic.py        # Vector Store Integration
+-- safety/
|   +-- validator.py       # Input Validation & Sanitization
|   +-- sandbox.py         # Docker-based Execution Env
+-- evaluation/
    +-- benchmarks.py      # Test Scenarios
    +-- metrics.py         # Success Rate, Costs, Latency
\end{lstlisting}

\subsection{Kernimplementierung: Agent-Controller}\index{Agent Controller!Implementierung}

Listing \ref{lst:agent-loop} zeigt eine minimale agentische Schleife\index{Agent Loop} mit Planung\index{Planung!Implementierung}, Werkzeugausführung und Reflexion\index{Reflexion}. Die vollständige Implementierung umfasst zusätzlich Fehlerbehandlung\index{Error-Handling!Agent}, Timeouts, Schrittbegrenzungen und strukturiertes Logging\index{Agent Logging}.

\begin{lstlisting}[style=Python, caption={Minimaler agentischer Loop für SE-Aufgaben (Beispiel-Listing)}, label={lst:agent-loop}]
from typing import Dict, Any

class Toolset:
  def run_tests(self) -> str:
    return "tests: 103 passed, 2 failed"

  def format_code(self, diff: str) -> str:
    return "formatted diff applied"

class Agent:
  def __init__(self, tools: Toolset):
    self.tools = tools
    self.memory = []  # episodic traces

  def plan(self, goal: str) -> str:
    return f"Plan: run tests -> fix failures -> re-run -> format -> commit ({goal})"

  def act(self, step: str) -> str:
    if "run tests" in step:
      return self.tools.run_tests()
    if "format" in step:
      return self.tools.format_code(diff="...")
    return "noop"

  def reflect(self, observation: str) -> str:
    if "failed" in observation:
      return "Next: inspect failing tests and patch code"
    return "Next: finalize and commit"

  def run(self, goal: str) -> Dict[str, Any]:
    plan = self.plan(goal)
    self.memory.append({"plan": plan})
    obs = self.act("run tests")
    self.memory.append({"obs": obs})
    next_step = self.reflect(obs)
    self.memory.append({"reflect": next_step})
    obs2 = self.act("format")
    self.memory.append({"obs": obs2})
    return {"status": "done", "trace": self.memory}

agent = Agent(Toolset())
result = agent.run(goal="increase reliability of module X")
print(result["status"])
\end{lstlisting}

Listing~\ref{lst:ts-tool-calling} zeigt ergänzend die Implementierung eines Tool-Aufrufs in TypeScript.

% ------------------------------------------------------------
% Zweites Listing: TypeScript Tool-Calling Stub
\begin{lstlisting}[style=TypeScript, caption={Werkzeugaufruf-Stub in TypeScript mit einfachem Funktionsschema (Beispiel-Listing)}, label={lst:ts-tool-calling}]
type ToolName = "run_tests" | "format_code" | "open_issue";

interface ToolCall {
  name: ToolName;
  args: Record<string, unknown>;
}

interface ToolResult {
  name: ToolName;
  ok: boolean;
  output: string;
}

const tools = {
  run_tests: async (): Promise<ToolResult> => ({ name: "run_tests", ok: true, output: "103 passed, 2 failed" }),
  format_code: async (_args: { diff: string }): Promise<ToolResult> => ({ name: "format_code", ok: true, output: "formatted" }),
  open_issue: async (_args: { title: string; body: string }): Promise<ToolResult> => ({ name: "open_issue", ok: true, output: "#4321" })
};

async function dispatch(call: ToolCall): Promise<ToolResult> {
  switch (call.name) {
    case "run_tests":
      return tools.run_tests();
    case "format_code":
      return tools.format_code(call.args as { diff: string });
    case "open_issue":
      return tools.open_issue(call.args as { title: string; body: string });
  }
}

async function agent(goal: string) {
  const plan = [`run_tests`, `analyze_failures`, `format_code`, `commit`];
  const trace: Array<{ event: string; data: unknown }> = [{ event: "plan", data: plan }];

  const res1 = await dispatch({ name: "run_tests", args: {} });
  trace.push({ event: "tool_result", data: res1 });

  if (res1.output.includes("failed")) {
    // Simple reflection -> open an issue with details
    const res2 = await dispatch({ name: "open_issue", args: { title: `Test failures for ${goal}`, body: res1.output } });
    trace.push({ event: "tool_result", data: res2 });
  }

  const res3 = await dispatch({ name: "format_code", args: { diff: "..." } });
  trace.push({ event: "tool_result", data: res3 });
  return { status: "done", trace };
}

agent("increase reliability of module X").then(r => console.log(r.status));
\end{lstlisting}

\section{Experimentelles Setup}

Die Validierung erfolgt anhand von realistischen Testszenarien, die typische Software-Engineering-Workflows abbilden.

\subsection{Testumgebung}

\begin{sloppypar}
\begin{description}
  \item[Hardware:] MacBook Pro M2, 16\,GB RAM, 512\,GB SSD
  \item[Betriebssystem:] macOS 14.x, Docker Desktop 4.25
  \item[LLM-Modelle:] GPT-4-Turbo (0125), Claude 3.5 Sonnet, mit Temperatur 0.2 für Reproduzierbarkeit
  \item[Testdaten:] 25 repräsentative Python-Projekte (5\,k--50\,k LOC), synthetische Bugs, real-world Issues
\end{description}
\end{sloppypar}

\subsection{Benchmark-Szenarien}

Drei Haupt-Szenarien wurden evaluiert (vgl. Kapitel~\ref{ch:konzept}):

\begin{enumerate}
  \item \textbf{Automated Refactoring} (10 Tasks): Extract-Function, Rename-Variable, Simplify-Conditional
  \item \textbf{Test Failure Diagnosis} (8 Tasks): Debug failing tests, fix assertions, update mocks
  \item \textbf{Lint Error Resolution} (7 Tasks): Fix style issues, type errors, unused imports
\end{enumerate}

Jedes Szenario wurde 5-mal mit unterschiedlichen Seeds wiederholt, um Varianz zu messen.

\subsection{Beispiel: Test Failure Diagnosis — Schritt-für-Schritt}

Im Folgenden wird das Szenario \enquote{Test Failure Diagnosis} detailliert beschrieben, um den praktischen Ablauf und die typischen Artefakte (Tool-Calls, Logs, Entscheidungen) zu veranschaulichen.

1) Initialer Zustand: Test-Runner meldet \texttt{3 failed} in \texttt{auth/test\_login.py}. Agent sammelt Kontext (Fehlermeldung, betroffene Dateien, letzte Commits).

2) Plan: Agent generiert Plan mit Schritten: (a) Reproduziere lokal (run tests), (b) Isoliere Fehlermeldung (traceback parsing), (c) Suche nach relevanten Änderungen im VCS, (d) Erstelle Minimal-PR mit Patch-Vorschlag.

3) Act: Tool-Calls (Beispiel):
\begin{itemize}
  \item \texttt{run\_tests(module="auth")} \textrightarrow\ \texttt{"3 failed, 97 passed"}
  \item \texttt{run\_linter(path="auth/")} \textrightarrow\ Tool-Result: Hinweise auf Stil, aber keine direkten Fehler
  \item \texttt{search\_in\_repo(query="login", range="last-5-commits")} \textrightarrow\ Treffer: Commit 123abc \texttt{Refactor: auth flow}
\end{itemize}

4) Observation: Agent parst Traceback, findet NullPointer-ähnlichen Fehler in Hilfsfunktion, die neu extrahiert wurde. Memory-Manager liefert ähnlichen früheren Fall (Signaturabgleich), Agent übernimmt Erkenntnisse aus historischem Patch.

5) Reflection 
\begin{sloppypar}
\begin{itemize}
  \item Agent generiert Patch-Vorschlag (kleine Scope-Korrektur im Helper), erstellt Diff und führt Tests erneut in isolierter Sandbox aus.
  \item Tests grün $\Rightarrow$ Agent eröffnet PR-Entwurf und notiert Review-Kommentare; bei Teil-Erfolg: Human-Approval-Gate vor Merge.
\end{itemize}
\end{sloppypar}

Beispiel-Trace (gekürzt):
\begin{lstlisting}[breaklines=true, basicstyle=\ttfamily\small]
[Agent] run_tests -> 3 failed (auth/test_login.py::test_login)
[Agent] search_in_repo -> found commit 123abc (Refactor auth flow)
[Agent] generate_patch -> diff created: modify auth/helpers.py
[Agent] run_tests (sandbox) -> 100 passed
[Agent] open_pr -> PR #42 (Draft)
\end{lstlisting}

Das Beispiel zeigt, wie Tool-Adapter, Memory und Safety-Layer (Sandbox, Human-Approval) zusammenwirken, um fehlerhafte Änderungen sicher zu erkennen und zu beheben. Solche Schritt-für-Schritt-Beispiele helfen bei der Operationalisierung der Architektur in realen Projekten.

\section{Ergebnisse}

Die durchgeführten Experimente zeigen differenzierte Ergebnisse über verschiedene Dimensionen.

\subsection{Erfolgsraten nach Aufgabentyp}

Tabelle \ref{tab:benchmark-results} zeigt detaillierte Ergebnisse für ausgewählte Tasks:

\begin{sloppypar}
\begin{itemize}
  \item \textbf{Refactoring:} \pct{73} Erfolgsrate (11/15 Tasks erfolgreich)
  \item \textbf{Test-Fixing:} \pct{62} Erfolgsrate (5/8 Tasks erfolgreich)
  \item \textbf{Lint-Resolution:} \pct{86} Erfolgsrate (6/7 Tasks erfolgreich)
\end{itemize}
\end{sloppypar}

Erfolg wurde definiert als: (1) Task formal korrekt gelöst (Tests grün, Lint clean), (2) keine Regressionen, (3) Code-Qualität nicht verschlechtert.

\subsection{Effizienzmetriken}

Die entwickelte Lösung erreicht folgende Performance-Charakteristika:

\begin{sloppypar}
\begin{description}
  \item[Token-Verbrauch:] Durchschnittlich 8.4\,k Tokens pro erfolgreichem Task (Range: 2\,k--25\,k). Kontextoptimierung reduzierte Verbrauch um \pct{40} gegenüber naivem Ansatz.
  
  \item[Laufzeit:] Median 12.3\,Sekunden pro Task (ohne Tool-Execution). Mit Test-Runs: 45\,s--180\,s je nach Testsuite-Größe.
  
  \item[Tool-Calls:] Durchschnittlich 4.2\,Werkzeugaufrufe pro Task. Reflexion reduzierte fehlerhafte Aufrufe um \pct{28}.
  
  \item[Kosten:] Geschätzt \$\,0.08 pro Task bei GPT-4-Pricing (Jan 2024). Claude war \pct{35} günstiger bei vergleichbarer Qualität.
\end{description}
\end{sloppypar}

\subsection{Qualitätsmetriken}

\textquote[\cite{yang2024sweagent}]{Die praktische Implementierung agentischer Systeme erfordert sorgfältige Planung und umfassende Tests, um Zuverlässigkeit in produktiven Umgebungen zu gewährleisten.}

Code-Qualität wurde über mehrere Dimensionen gemessen:

\begin{sloppypar}
\begin{itemize}
  \item \textbf{Test-Pass-Rate:} \pct{98} (nur 2 von 103 Tests regressierten nach Agent-Edits)
  \item \textbf{Lint-Score:} Durchschnittlich +12 Punkte Verbesserung nach Lint-Fixes
  \item \textbf{Complexity:} Cyclomatic Complexity blieb unverändert oder verbesserte sich (Extract-Function Tasks)
  \item \textbf{Review-Akzeptanz:} Manuelles Review ergab \pct{81} Akzeptanzrate („would merge“)
  \item \textbf{Review-Akzeptanz:} Manuelles Review ergab \pct{81} Akzeptanzrate (\enquote{would merge})
\end{itemize}
\end{sloppypar}

\subsection{Robustheit und Fehlerbehandlung}

Error-Cases wurden systematisch getestet:

\begin{sloppypar}
\begin{itemize}
  \item \textbf{Werkzeugfehler:} Agent erholte sich in \pct{67} der Fälle durch Wiederholung oder Alternativstrategie
  \item \textbf{Malformed Outputs:} JSON-Parsing-Fehler wurden durch Wiederholungen mit Schema-Validierung in \pct{89} behoben
  \item \textbf{Timeouts:} Graceful Degradation bei max-steps Limit (definiert als \pct{15} Partial-Success)
\end{itemize}
\end{sloppypar}

\section{Vergleich mit existierenden Ansätzen}

Die entwickelte Lösung wurde mit Baselines verglichen:

\begin{table}[ht]
\centering
\caption{Vergleich mit Baseline-Systemen (auf gleichem Benchmark-Set)}
\label{tab:comparison}
\begin{tabularx}{\textwidth}{Xcccc}
\toprule
\cellcolor{headercolor}\textcolor{white}{\textbf{Ansatz}} & \cellcolor{headercolor}\textcolor{white}{\textbf{Success}} & \cellcolor{headercolor}\textcolor{white}{\textbf{Token}} & \cellcolor{headercolor}\textcolor{white}{\textbf{Zeit (s)}} & \cellcolor{headercolor}\textcolor{white}{\textbf{Kosten}} \\
\midrule
Naive Prompting & \pct{42} & 12.3\,k & 8.5 & \$\,0.12 \\
\rowcolor{rowcolorgray}
ReAct (Baseline) & \pct{58} & 10.1\,k & 15.2 & \$\,0.10 \\
Unsere Architektur & \pct{73} & 8.4\,k & 12.3 & \$\,0.08 \\
\rowcolor{rowcolorgray}
+ Reflexion & \pct{73} & 8.9\,k & 14.1 & \$\,0.09 \\
\bottomrule
\end{tabularx}
\end{table}

Kernverbesserungen gegenüber Baselines:

\begin{sloppypar}
\begin{itemize}
  \item \textbf{+31\% Erfolgsrate} vs. Naives Prompting durch strukturierte Werkzeugorchestrierung
  \item \textbf{+15\% Erfolgsrate} vs. Standard-ReAct durch optimiertes Kontextmanagement
  \item \textbf{-17\% Token-Kosten} durch intelligentes Pruning und Zusammenfassung
  \item \textbf{Robustere Fehlerbehebung} durch Reflexions-Mechanismen
\end{itemize}
\end{sloppypar}

Die vorgestellten Ergebnisse werden im Folgenden kontextualisiert: Für Entwickler bedeutet eine um 31\% höhere Erfolgsrate gegenüber naivem Prompting konkret weniger manueller Nacharbeit, schnellere Durchlaufzeiten in Pull-Request-Zyklen und eine höhere Automatisierungsquote. Für Betreiber von CI/CD-Infrastrukturen bedeuten niedrigere Token-Kosten und reduzierte Fehlerraten geringere Betriebskosten. Die praktische Perspektive ist wichtig, um Entscheidungsträgern in Unternehmen handfeste Argumente für den Einsatz agentischer Lösungen zu liefern.

\section{Validierung und Verifikation}

Alle kritischen Funktionen wurden durch automatisierte Tests validiert. Die Testabdeckung beträgt \pct{87} (Zeilen-Coverage).

\subsection{Unit-Tests}

\begin{sloppypar}
\begin{itemize}
  \item Tool-Adapter: 45 Tests, \pct{95} Abdeckung
  \item Policy-Logic: 32 Tests, \pct{89} Abdeckung
  \item Sicherheitsschicht: 28 Tests, \pct{92} Abdeckung
\end{itemize}
\end{sloppypar}

\subsection{Integrationstests}

End-to-End-Tests für alle 3 Benchmark-Szenarien. Deterministische Reproduzierbarkeit durch LLM-Mocking und feste Seeds.

Praktische Erkenntnisse: Automatisierte Tests sollten sowohl synthetische als auch realistische Repositories umfassen; Mocking allein reicht nicht aus, da Tests in realen Umgebungen unerwartete Randfälle aufdecken. Darüber hinaus ist kontinuierliche Überwachung unabdingbar, um Drift in LLM-Verhalten oder Änderungen in abhängigen Tools frühzeitig zu erkennen.

\subsection{Sicherheits-Audits}

\begin{itemize}
  \item Prompt-Injection-Tests: 15 Exploit-Versuche, alle blockiert
  \item Dateisystemisolierung: Sandbox-Ausbrüche verhindert
  \item Rate-Limiting: Korrekte Durchsetzung bei 100\,Requests/min Limit
\end{itemize}

% Benchmark-Ergebnisse mit longtable
\begin{table}[ht]
\centering
\caption{Detaillierte Benchmark-Ergebnisse: Agentische SE-Tasks mit Evaluationsmetriken}
\label{tab:benchmark-results}
\begin{tabularx}{\textwidth}{Xccccc}
\toprule
\cellcolor{headercolor}\textcolor{white}{\textbf{Aufgabe}} & \cellcolor{headercolor}\textcolor{white}{\textbf{Success}} & \cellcolor{headercolor}\textcolor{white}{\textbf{Token}} & \cellcolor{headercolor}\textcolor{white}{\textbf{Zeit (s)}} & \cellcolor{headercolor}\textcolor{white}{\textbf{Tools}} & \cellcolor{headercolor}\textcolor{white}{\textbf{Kosten}} \\
\midrule
Extract Function & OK & 7.2\,k & 18.5 & 4 & \$\,0.07 \\
\rowcolor{rowcolorgray}
Fix Lint Errors & OK & 3.1\,k & 8.2 & 3 & \$\,0.03 \\
Debug Test Failure & OK & 12.5\,k & 45.3 & 6 & \$\,0.12 \\
\rowcolor{rowcolorgray}
Format Code & OK & 2.8\,k & 5.1 & 2 & \$\,0.03 \\
Rename Variable & OK & 5.4\,k & 12.8 & 3 & \$\,0.05 \\
\rowcolor{rowcolorgray}
Remove Deadcode & OK & 8.9\,k & 22.1 & 5 & \$\,0.09 \\
Update Mocks & FAIL & 9.7\,k & 38.2 & 7 & \$\,0.10 \\
\rowcolor{rowcolorgray}
Simplify Conditional & OK & 6.3\,k & 15.7 & 4 & \$\,0.06 \\
Generate Docstrings & OK & 4.5\,k & 9.3 & 2 & \$\,0.04 \\
\rowcolor{rowcolorgray}
Fix Type Errors & OK & 11.2\,k & 28.4 & 5 & \$\,0.11 \\
\bottomrule
\multicolumn{6}{l}{\small \textit{Durchschnitt (erfolgreiche Tasks):} 7.1\,k Tokens, 16.5\,s, 3.8\,Tools, \$\,0.07} \\
\end{tabularx}
\end{table}

Tabelle~\ref{tab:eval-metrics} fasst die verwendeten Evaluationsmetriken zusammen.

% Zusätzliche Beispiel-Tabelle für Evaluationsmetriken
\begin{table}[ht]
\centering
\caption{Evaluationsmetriken für agentische SE-Workflows (Beispieltabelle)}
\label{tab:eval-metrics}
\begin{tabularx}{\textwidth}{lX}
\toprule
\cellcolor{headercolor}\textcolor{white}{\textbf{Kategorie}} & \cellcolor{headercolor}\textcolor{white}{\textbf{Metriken / Beschreibung}} \\
\midrule
Qualität & Task-Success-Rate, Patch-Korrektheit (Tests/Lint), Review-Akzeptanz, Regressionen (\#) \\
\rowcolor{rowcolorgray}
Kosten & Token-/API-Kosten (EUR), Werkzeugaufrufkosten, Rechenzeit \\
Latenz & End-to-End-Laufzeit (s), Tool-Roundtrips (\#), Wartezeit auf CI \\
\rowcolor{rowcolorgray}
Sicherheit & Policy-Verstöße (\#), Risk Flags, sand-boxed I/O, PII-Leaks (\#) \\
Nachvollziehbarkeit & Trace-Länge (\#Events), Artefakte (Patches, Logs), Reproduzierbarkeit (Seeds) \\
\bottomrule
\end{tabularx}
\end{table}
